\documentclass[11pt]{article}
\usepackage[margin=1in]{geometry}
\usepackage{amsmath,amssymb,amsthm,mathtools}
\usepackage{graphicx}
\usepackage{microtype}
\usepackage[hidelinks]{hyperref}
\usepackage[nameinlink,noabbrev]{cleveref}

\newtheorem{theorem}{Theorem}
\newtheorem{proposition}[theorem]{Proposition}
\newtheorem{lemma}[theorem]{Lemma}
\newtheorem{corollary}[theorem]{Corollary}
\theoremstyle{remark}
\newtheorem{remark}[theorem]{Remark}

\newcommand{\R}{\mathbb R}
\newcommand{\g}{\mathfrak g}
\newcommand{\h}{\mathfrak h}
\newcommand{\dd}{\,\mathrm d}
\newcommand{\abs}[1]{\lvert #1\rvert}
\newcommand{\norm}[1]{\lVert #1\rVert}
\newcommand{\Vq}[3]{\mathcal V_{#1,#2}(#3)}
\newcommand{\Adiff}[3]{A_{#1}^{#3}(#2)}

\title{Endpoint \(L^1\) Vertical Inequalities in Nilpotency Step Two\\
Verification, Strengthening, and Literature Status}
\author{}
\date{June 18, 2026}

\begin{document}
\maketitle

\begin{abstract}
We verify the argument in the supplied solution packet for Ryoo's endpoint
vertical-versus-horizontal question.  The proof is correct after making two standard
points explicit: the elementary Heisenberg subgroup is closed, and Weil's quotient
integration formula is applied to that closed unimodular subgroup.  We also strengthen
its scope.  The Carnot assumption is unnecessary: every connected, simply connected,
step-two nilpotent Lie group equipped with an arbitrary bracket-generating left-invariant
Carnot--Carath\'eodory structure satisfies the endpoint inequality with exponent \(q=4\)
in every commutator direction.  A quantitative constant is given by a natural factorization
quasinorm of the central vector.  Hence Ryoo's Question 38 is completely resolved in
nilpotency step two, though the optimal exponent can be smaller than four.  We did not
find a published or preprint resolution, or a counterexample, for the general higher-step
question, and the higher-step problem remains open on the evidence located in the
literature search described below.
\end{abstract}

\section*{Packet status}

This is a strengthened partial packet for Ryoo's endpoint \(L^1\)
vertical-versus-horizontal problem.  It is a complete positive answer in nilpotency
step two, including non-Carnot step-two groups equipped with an arbitrary
bracket-generating left-invariant Carnot--Carath\'eodory structure.  It is not a
full solution of the higher-step endpoint problem.

\section{The endpoint question}
Let \(G\) be a connected, simply connected nilpotent Lie group, let \(W\subseteq\g\)
be a bracket-generating subspace with a Euclidean norm, and let \(\nabla_W f\) denote
the corresponding left-invariant horizontal gradient.  If \(v\) lies in the last nonzero
term of the lower central series and the nilpotency step is \(s\), Ryoo asks for which
\(q\geq 1\) one has
\begin{equation}\label{eq:Ryoo}
 \left(\int_0^\infty
 \left(\frac{\displaystyle\int_G
 \abs{f(g\exp(tv))-f(g)}\dd g}{t^{1/s}}\right)^q
 \frac{\dd t}{t}\right)^{1/q}
 \lesssim_{G,q,v}\int_G\abs{\nabla_W f(g)}\dd g
\end{equation}
for every \(f\in C_c^\infty(G)\).  The usual essential-supremum interpretation is used
when \(q=\infty\).

\begin{center}
\includegraphics[width=0.94\linewidth]{figures/open_problem_crop.png}
\end{center}

The supporting Heisenberg endpoint theorem used in the proof is the
three-dimensional \(q=4\) result of Naor and Young:

\begin{center}
\includegraphics[width=0.90\linewidth]{figures/supporting_h3_theorem_crop.png}
\end{center}

In the rest of the note, \(G\) has nilpotency step two.  Put
\[
 Z=[\g,\g].
\]
Then \(Z\) is central.  Since \(W\) Lie-generates \(\g\) and the group has step two,
\begin{equation}\label{eq:bracket-spans}
 \g=W+[W,W],\qquad Z=[W,W].
\end{equation}
For \(z\in Z\), \(t>0\), and \(f\in L^1(G)\), set
\[
 A_z^f(t)=\int_G\abs{f(g\exp(tz))-f(g)}\dd g
\]
and, for \(1\le q<\infty\),
\[
 \mathcal V_{q,z}(f)=
 \left(\int_0^\infty\left(\frac{A_z^f(t)}{\sqrt t}\right)^q
 \frac{\dd t}{t}\right)^{1/q},
 \qquad
 \mathcal V_{\infty,z}(f)=\sup_{t>0}\frac{A_z^f(t)}{\sqrt t}.
\]

A useful quantitative measure of the cost of writing \(z\) as a sum of horizontal
commutators is
\begin{equation}\label{eq:kappa}
 \kappa_W(z)=\inf\left\{
 \sum_{j=1}^m\sqrt{\abs{a_j}\abs{b_j}}:
 z=\sum_{j=1}^m[a_j,b_j],\quad a_j,b_j\in W
 \right\}.
\end{equation}
By \eqref{eq:bracket-spans}, \(\kappa_W(z)<\infty\).  Notice the natural homogeneity
\(\kappa_W(cz)=\sqrt{\abs c}\,\kappa_W(z)\).

\section{Complete step-two theorem}

\begin{theorem}[Endpoint inequality in every step-two group]\label{thm:main}
There is a universal constant \(C_4<\infty\) such that, for every connected, simply
connected, step-two nilpotent Lie group \(G\), every bracket-generating horizontal
space \(W\), every \(z\in[\g,\g]\), and every \(f\in C_c^\infty(G)\),
\begin{equation}\label{eq:q4}
 \mathcal V_{4,z}(f)
 \le C_4\,\kappa_W(z)\int_G\abs{\nabla_W f(g)}\dd g.
\end{equation}
Moreover, for every \(q\in[4,\infty]\),
\begin{equation}\label{eq:qge4}
 \mathcal V_{q,z}(f)
 \le C_q\,\kappa_W(z)\int_G\abs{\nabla_W f(g)}\dd g,
\end{equation}
where \(C_q\) is universal apart from its dependence on \(q\).
Consequently, Ryoo's Question 38 has a positive answer for every step-two nilpotent
Lie group, for every allowed final-layer vector, and for every \(q\ge4\).
\end{theorem}

The proof uses the sharp \(q=4\) theorem of Naor and Young on the three-dimensional
Heisenberg group.

\begin{lemma}[Transfer from an elementary Heisenberg subgroup]\label{lem:H3-transfer}
Let \(a,b\in W\) and \(z=[a,b]\ne0\).  Then
\begin{equation}\label{eq:elementary-unbalanced}
 \mathcal V_{4,z}(f)
 \le C_{\mathrm{NY}}\int_G\bigl(\abs{af(g)}+\abs{bf(g)}\bigr)\dd g
\end{equation}
for every \(f\in C_c^\infty(G)\), where \(C_{\mathrm{NY}}\) is the universal constant
in the three-dimensional Heisenberg theorem.  In particular,
\begin{equation}\label{eq:elementary-balanced}
 \mathcal V_{4,z}(f)
 \le 2C_{\mathrm{NY}}\sqrt{\abs a\abs b}
 \int_G\abs{\nabla_W f(g)}\dd g.
\end{equation}
\end{lemma}

\begin{proof}
The vectors \(a,b,z\) are linearly independent.  Indeed, if the images of \(a\) and
\(b\) in \(\g/Z\) were dependent, then \([a,b]=0\).  Hence
\[
 \h=\operatorname{span}\{a,b,z\}
\]
is a three-dimensional Heisenberg Lie algebra with \([a,b]=z\).  Since the exponential
map of a simply connected nilpotent group is a diffeomorphism, the connected subgroup
\(H=\exp(\h)\) is closed: it is the image under a homeomorphism of the closed linear
subspace \(\h\subseteq\g\).

Both \(G\) and \(H\) are unimodular.  Weil's quotient integration formula therefore
gives a \(G\)-invariant measure on \(G/H\) such that
\begin{equation}\label{eq:Weil}
 \int_G\Phi(g)\dd g
 =\int_{G/H}\int_H\Phi(xu)\dd u\,\dd\bar x
\end{equation}
for compactly supported measurable \(\Phi\).  For a coset representative \(x\), define
\(F_x(u)=f(xu)\).  The Lie algebra isomorphism from the standard Heisenberg algebra
that sends its standard basis \((X,Y,Z_0)\) to \((a,b,z)\) identifies the Naor--Young
inequality with
\begin{equation}\label{eq:coset-NY}
 \left(\int_0^\infty
 \left(\frac{\displaystyle\int_H
 \abs{f(xu\exp(tz))-f(xu)}\dd u}{\sqrt t}\right)^4
 \frac{\dd t}{t}\right)^{1/4}
 \le C_{\mathrm{NY}}
 \int_H\bigl(\abs{af(xu)}+\abs{bf(xu)}\bigr)\dd u.
\end{equation}
A scalar change in the normalization of Haar measure multiplies both sides of
\eqref{eq:coset-NY} by the same factor.  Also, the left-invariant derivatives of
\(F_x\) in the directions \(a,b\) are exactly the ambient derivatives \(af,bf\) at
\(xu\).

Let
\[
 B_x(t)=\int_H\abs{f(xu\exp(tz))-f(xu)}\dd u.
\]
By \eqref{eq:Weil}, \(A_z^f(t)=\int_{G/H}B_x(t)\dd\bar x\).  Minkowski's integral
inequality in \(L^4((0,\infty),\dd t/t)\), followed by \eqref{eq:coset-NY} and
\eqref{eq:Weil}, gives \eqref{eq:elementary-unbalanced}.

For \(\lambda>0\), the same central vector has the factorization
\[
 z=[\lambda a,\lambda^{-1}b].
\]
Applying \eqref{eq:elementary-unbalanced} to this pair and using
\(\abs{cf}\le\abs c\,\abs{\nabla_W f}\) for \(c\in W\), we obtain
\[
 \mathcal V_{4,z}(f)
 \le C_{\mathrm{NY}}
 \bigl(\lambda\abs a+\lambda^{-1}\abs b\bigr)
 \int_G\abs{\nabla_Wf}\dd g.
\]
Minimization in \(\lambda\) yields \eqref{eq:elementary-balanced}.
\end{proof}

\begin{proof}[Proof of \cref{thm:main}]
Fix a representation \(z=\sum_{j=1}^m[a_j,b_j]\).  Set
\(z_j=[a_j,b_j]\) and \(w_j=z_1+\cdots+z_j\), with \(w_0=0\).  Since \(Z\) is
central,
\[
 \exp(tz)=\prod_{j=1}^m\exp(tz_j).
\]
For each \(t>0\), telescoping and right invariance of Haar measure give
\begin{align*}
 A_z^f(t)
 &\le\sum_{j=1}^m
 \int_G\abs{f(g\exp(tw_j))-f(g\exp(tw_{j-1}))}\dd g\\
 &=\sum_{j=1}^m A_{z_j}^f(t).
\end{align*}
Minkowski's inequality and \cref{lem:H3-transfer} imply
\[
 \mathcal V_{4,z}(f)
 \le 2C_{\mathrm{NY}}
 \left(\sum_{j=1}^m\sqrt{\abs{a_j}\abs{b_j}}\right)
 \int_G\abs{\nabla_Wf}\dd g.
\]
Taking the infimum over all factorizations proves \eqref{eq:q4}.

For completeness, we record the deterministic endpoint.  In a step-two group,
\[
 \exp(ra)\exp(ub)\exp(-ra)\exp(-ub)=\exp(ru[a,b]).
\]
Thus \(\exp(t[a,b])\) can be joined to the identity by a horizontal commutator loop
of length at most \(4\sqrt{t\abs a\abs b}\), after optimizing over \(ru=t\).
Concatenating these loops and taking the infimum in \eqref{eq:kappa} gives
\begin{equation}\label{eq:distance}
 d_G(e,\exp(tz))\le4\sqrt t\,\kappa_W(z).
\end{equation}
If \(\gamma\) is a horizontal path from \(e\) to \(\exp(tz)\), differentiation of
\(f(g\gamma(\cdot))\), followed by Fubini and right invariance of Haar measure, yields
\[
 A_z^f(t)\le \operatorname{length}(\gamma)
 \int_G\abs{\nabla_Wf}\dd g.
\]
Together with \eqref{eq:distance}, this proves
\begin{equation}\label{eq:qinf}
 \mathcal V_{\infty,z}(f)
 \le4\kappa_W(z)\int_G\abs{\nabla_Wf}\dd g.
\end{equation}
Interpolation of the scalar function \(t\mapsto A_z^f(t)/\sqrt t\) between
\(L^4(\dd t/t)\) and \(L^\infty(\dd t/t)\) proves \eqref{eq:qge4}.
\end{proof}

\begin{corollary}[The supplied Carnot statement]\label{cor:Carnot}
If \(G\) is a step-two Carnot group with stratification \(\g=V_1\oplus V_2\), then
\cref{thm:main} applies with \(W=V_1\) and every \(z\in V_2\).  In particular, the
main theorem in the supplied packet is correct.
\end{corollary}

\begin{corollary}[Discrete consequence]\label{cor:discrete}
Combining \cref{thm:main} with Ryoo's discretization argument gives exponent \(q=4\)
in the corresponding discrete endpoint inequality for finitely generated step-two
nilpotent groups, and hence for the step-two cases covered after passing to finite-index
and finite-torsion modifications in Ryoo's framework.
\end{corollary}

\section{A sufficient mechanism for exponent two}
The exponent four is not always optimal.  Naor and Young proved exponent two for
\(\mathbb H^{2k+1}\) when \(k\ge2\).  The same coset-transfer argument gives the
following useful sufficient criterion.

\begin{proposition}[Higher-dimensional Heisenberg transfer]\label{prop:q2}
Suppose that a central vector \(z\in[\g,\g]\) generates the center of a closed connected
subgroup \(H\le G\) whose Lie algebra is isomorphic to \(\mathfrak h^{2k+1}\) for some
\(k\ge2\), and suppose that the horizontal layer of \(H\) is contained in \(W\).  Then
\[
 \mathcal V_{2,z}(f)\lesssim_{H,W}
 \int_G\abs{\nabla_Wf}\dd g.
\]
Consequently the same estimate holds for every \(q\ge2\).  More generally, if a vector
\(v\) is a finite sum of central vectors satisfying this criterion, then the conclusion
holds for \(v\) by central telescoping.
\end{proposition}

\begin{proof}
Disintegrate Haar measure over \(G/H\), apply the higher-dimensional Heisenberg
\(q=2\) theorem on each coset, and use Minkowski exactly as in
\cref{lem:H3-transfer}.  Central telescoping handles finite sums, and interpolation with
\eqref{eq:qinf} gives all larger exponents.
\end{proof}

\section{Verification of the supplied proof}
The supplied argument has the right structure and no fatal gap.  The following are the
points that require, or benefit from, explicit wording.

\begin{enumerate}
\item The subgroup generated by \(a,b,[a,b]\) is closed.  In a simply connected
nilpotent group this follows immediately from the fact that the exponential map is a
diffeomorphism.
\item Weil disintegration is applicable because the subgroup is closed and both groups
are unimodular.
\item For \(F_x(u)=f(xu)\), the subgroup left-invariant derivatives are the ambient
left-invariant derivatives at \(xu\).
\item Haar normalization under the Lie-group identification with \(\mathbb H^3\) is
harmless because both sides of the endpoint inequality are homogeneous of degree one
in Haar measure.
\item The final telescoping identity uses centrality and right invariance of Haar measure.
Nilpotent groups are unimodular, so this is valid.
\end{enumerate}

The principal strengthening over the packet is that no grading or Carnot dilation is used
in the proof.  The only algebraic inputs are step two and the fact that the chosen horizontal
space bracket-generates the Lie algebra.

\section{Literature status as of June 18, 2026}
Ryoo's revised paper states the endpoint problem as Questions 38 and 42 and records only
the exact Heisenberg answers: exponent range \([4,\infty)\) in \(\mathbb H^3\), and
\([2,\infty)\) in \(\mathbb H^{2k+1}\), \(k\ge2\).  It then formulates the existence of a
finite exponent in every nilpotent group as a conjecture.  The article was accepted in
March 2026 and made available online in April 2026; the endpoint questions remain in the
revised version.

We searched exact phrases from Questions 38 and 42 and Conjecture 43, the arXiv record
and papers citing it, combinations of ``endpoint'', ``vertical perimeter'', ``nilpotent'',
``Carnot'', and ``\(L^1\)'', and recent 2025--2026 Carnot-group BV and nonlocal Sobolev
papers.  The search located the Heisenberg theorems, the qualitative theorem that general
nonabelian nilpotent groups do not embed bi-Lipschitzly into \(L^1\), and recent isotropic
or horizontal nonlocal characterizations of BV functions.  None supplies the one-directional,
last-layer, finite-\(q\) estimate in \eqref{eq:Ryoo}, and no counterexample to it was located.
Thus, on the evidence found, there is no known full answer beyond the Heisenberg cases
and the complete step-two upper bound proved above.

This is necessarily a literature-search conclusion rather than a theorem: absence from the
sources searched cannot certify absolute nonexistence of an unpublished or differently
phrased result.

\section{Why the elementary proof does not iterate to higher step}
It is useful to isolate the precise obstruction.  Suppose \(v=[a,w]\) lies in the top
layer of a step-\(s\) group, where \(a\in W\) and \(w\) has homogeneous degree \(s-1\).
The subgroup generated by \(a,w,v\) is three-dimensional Heisenberg, because \(v\) is
central.  Applying the Heisenberg theorem there, however, introduces the derivative
\(wf\), which is not controlled by the first-order horizontal variation of \(f\).

At the level of finite differences one has the exact commutator path estimate
\begin{equation}\label{eq:heuristic}
 A_v^f(ru)\le2A_a^f(r)+2A_w^f(u).
\end{equation}
Writing \(t=ru\), choosing \(r=\lambda t^{1/s}\), and formally normalizing the
\(w\)-term at its degree gives
\[
 \frac{A_v^f(t)}{t^{1/s}}
 \lesssim \lambda\int_G\abs{\nabla_Wf}
 +\lambda^{-1/(s-1)}
 \frac{A_w^f(\lambda^{-1}t^{(s-1)/s})}
 {(\lambda^{-1}t^{(s-1)/s})^{1/(s-1)}}.
\]
This suggests a distribution-function induction and even predicts finite candidate
exponents.  The induction does not close: \(w\) is generally not central in \(G\), while
it becomes central only after quotienting by the top layer.  A section of that quotient
introduces a top-layer cocycle into the horizontal derivatives, precisely the quantity one
is trying to estimate.  In the Engel group this obstruction is already present with
\(v=[X,[X,Y]]\) and \(w=[X,Y]\).

Other approaches examined include nested commutator words, quotient induction,
constant-normal sets in the Engel and free rank-two step-three groups, and recent
nonlocal BV characterizations.  They give useful pointwise or qualitative information,
but no scale-summable estimate and no counterexample.  In particular, the known
pathological constant-normal examples have top-central translation moduli that decay
faster than the critical \(t^{1/s}\) normalization in the model calculations, so they do not
by themselves disprove the conjecture.

\section{Conclusion}
The partial result is correct and can be upgraded to a complete theorem for all
step-two nilpotent Lie groups with arbitrary bracket-generating left-invariant
Carnot--Carath\'eodory structures.  The proof is short because every commutator vector is
a finite sum of centers of genuine horizontal three-dimensional Heisenberg subgroups.
No analogous reduction is available in higher step: the second generator of the elementary
Heisenberg subgroup is then a nonhorizontal lower-central direction.  We did not obtain a
rigorous higher-step proof or counterexample, and none was found in the literature search.

\begin{thebibliography}{99}

\bibitem{Ryoo}
S.-Y. Ryoo,
\emph{Quantitative nonembeddability of groups of polynomial growth into uniformly convex spaces},
J. Funct. Anal. (2026), article 111491; arXiv:2207.11305v2.

\bibitem{NY22}
A. Naor and R. Young,
\emph{Foliated corona decompositions},
Acta Math. \textbf{229} (2022), 55--200.

\bibitem{NY18}
A. Naor and R. Young,
\emph{Vertical perimeter versus horizontal perimeter},
Ann. of Math. (2) \textbf{188} (2018), 171--279.

\bibitem{EBGLD}
S. Eriksson-Bique, C. Gartland, E. Le Donne, L. Naples, and S. Nicolussi Golo,
\emph{Nilpotent groups and bi-Lipschitz embeddings into \(L^1\)},
Int. Math. Res. Not. IMRN (2023), no. 12, 10759--10797.

\bibitem{SCZ}
F. Serra Cassano and K. Zambanini,
\emph{Characterizations of Sobolev and BV functions on Carnot groups},
arXiv:2603.28996, 2026.

\bibitem{BLD}
C. Bellettini and E. Le Donne,
\emph{Sets with constant normal in Carnot groups: properties and examples},
Comment. Math. Helv. \textbf{96} (2021), 149--198.

\end{thebibliography}

\end{document}
